\section{Introduction}
\label{sec:intro}

Coding-agent leaderboards now rank pairs of an agent and a model rather than models alone
\citep{tbench2026board,aa2026coding,kapoor2026hal}, and many products wrap the same few frontier
models. For the people who build the wrapper, the harness, meaning the prompts, tools, context
management, team structure and runtime checks around the model, the practical question is what that
harness adds once the model is fixed. The published answers disagree. Multi-agent frameworks report
gains from role-playing teams, debate and voting
\citep{wu2024autogen,hong2024metagpt,qian2024chatdev,du2024debate,li2024moreagents}; a failure
analysis of such systems finds that their gains over a single agent are often small
\citep{cemri2025mast}; and practitioners argue both for a single agent with shared context
\citep{yan2025dontbuild} and for a lead that fans work out to parallel helpers
\citep{anthropic2025multiagent}.

We build The One, a product that runs teams of language-model agents on a person's computer or on
a Linux server. Between 28 September and 20 November 2026 we ran three experimental programs on it
to answer four questions:
\begin{enumerate}[label=\textbf{RQ\arabic*},leftmargin=3.2em]
  \item What does context management buy in cost, latency and accuracy?
  \item Do team structures beat one capable agent, and on which tasks?
  \item Can the harness steer the model's judgment at the points where runs fail?
  \item Where does the resulting configuration stand against the model vendor's own agent?
\end{enumerate}

\para{Contributions} The contributions are a record and a method more than a system.
\begin{itemize}[leftmargin=1.4em]
  \item \emph{A record of controlled experiments.} Over 120 batches on our own task suite (V1),
    315 Harbor jobs with 6,223 graded Terminal-Bench~4.0 trials (V8), and a Terminal-Bench~4.1 run
    under the Artificial Analysis (AA) protocol, with the rule, the branches and our own forecast
    written down before every batch ran. An index of every experiment is in \cref{app:index}.
  \item \emph{A held-out discipline with measured resolution.} Thirty Terminal-Bench tasks were never
    opened; every arm ran beside its base in the same hours; differences were tested by permutation
    within tasks; zeros were re-graded on an idle machine. The discipline caught most of our own
    false wins, and we state the smallest effect it can see (\cref{sec:method}).
  \item \emph{Measured effects, mostly negative.} The effect of context levers on cost and latency
    (\cref{sec:context}); a failure pattern shared by team structures and same-model second opinions
    (\cref{sec:structure}); and a rule for steering the model (a note changes the process, a fact
    changes the choice; \cref{sec:judgment}).
  \item \emph{A public adapter.} The Harbor adapter and the exact run command are public
    \citep{zhang2026adapter}; the configuration ties the vendor's agent on Terminal-Bench~4.1 with
    fewer input tokens (\cref{sec:tb41}).
\end{itemize}

\para{Summary of findings} Append-only prompts cut cost per run by a third to a half at equal
scores, and later context levers took off another 7 to 20\%; none of them moved accuracy. On tasks
that one context can hold, a single strong owner matched or beat every team we built, usually at a
fraction of the cost, because the failures that remain are single misreadings that every copy of
the same model shares. Text that advised the model changed what it did next but not what it chose;
facts and measurements placed in front of it did change the choice. The model itself set the
ceiling: moving from GPT-6.1 Sol to GPT-6.2 Sol changed which tasks passed more than any harness
lever did.

\para{Scope} We do not claim a new model, a result for tasks too big for one context, or a
leaderboard position. All runs were self-run on one machine, mostly through a subscription
sign-in, and costs are reconstructed at API list prices (\cref{sec:threats}).

The rest of the paper is organized as follows. \Cref{sec:related} reviews related work,
\cref{sec:system} describes the system and \cref{sec:method} the evaluation method.
\Cref{sec:context,sec:structure,sec:judgment,sec:person,sec:tb41} report results,
\cref{sec:discussion} explains why most levers failed, \cref{sec:threats} lists threats to
validity, and \cref{sec:conclusion} concludes.
